% Gabarit pandoc des corrigés publiés (brut et guidé). Contenu sans marque ni nom de site.
% Compilation : lualatex (latexmk -lualatex).
% Le Markdown reste la source unique : ne pas éditer ce fichier .tex à la main.
\documentclass[11pt,a4paper]{article}

% --- Moteur, polices (Latin Modern) et langue ---------------------------------
\usepackage{iftex}
\ifLuaTeX\else\errmessage{Ce document se compile avec lualatex}\fi
\usepackage{amsmath,amssymb,mathtools}
\usepackage{fontspec}
% mathrm=sym : \mathrm prend ses glyphes dans la police mathématique (ℓ, μ, Ω), comme KaTeX.
\usepackage[math-style=TeX,bold-style=TeX,mathrm=sym]{unicode-math}
% Repli pour les glyphes absents de Latin Modern (indices et exposants Unicode,
% flèches, filets des schémas en texte, symboles) : DejaVu, puis Latin Modern Math
% (ℓ, lettres mathématiques Unicode des légendes).
\directlua{
  luaotfload.add_fallback("repliTexte", {"DejaVuSerif:mode=harf;", "DejaVuSans:mode=harf;", "file:latinmodern-math.otf:mode=harf;"})
  luaotfload.add_fallback("repliMono", {"DejaVuSansMono:mode=harf;", "DejaVuSans:mode=harf;", "file:latinmodern-math.otf:mode=harf;"})
}
\setmainfont{Latin Modern Roman}[RawFeature={fallback=repliTexte}]
\setsansfont{Latin Modern Sans}[RawFeature={fallback=repliTexte}]
\setmonofont{Latin Modern Mono}[RawFeature={fallback=repliMono}]
\setmathfont{Latin Modern Math}
\usepackage[french]{babel}
\frenchsetup{SmallCapsFigTabCaptions=false}

% --- Mise en page --------------------------------------------------------------
\usepackage[a4paper,margin=2.2cm,headheight=14pt,footskip=1.2cm]{geometry}
\usepackage{microtype}
\setlength{\emergencystretch}{4em}
\tolerance=1500
\setlength{\parindent}{0pt}
\setlength{\parskip}{0.55\baselineskip plus 2pt minus 1pt}
\usepackage{xcolor}
% Gris nommés : babel-french rend « ! » actif dans le document.
\definecolor{gris2}{gray}{0.98}
\definecolor{gris3}{gray}{0.97}
\definecolor{gris5}{gray}{0.95}
\definecolor{gris35}{gray}{0.65}
\definecolor{gris40}{gray}{0.60}
\definecolor{gris45}{gray}{0.55}
\definecolor{gris50}{gray}{0.50}
\definecolor{gris55}{gray}{0.45}
\definecolor{gris60}{gray}{0.40}
\definecolor{gris65}{gray}{0.35}
\definecolor{gris70}{gray}{0.30}
\definecolor{gris80}{gray}{0.20}
\definecolor{gris85}{gray}{0.15}
\usepackage{graphicx}
\usepackage{float}
\usepackage{caption}
\captionsetup{font=small,labelformat=empty,justification=centering}
\usepackage{booktabs,longtable,array,multirow,calc}
\usepackage{siunitx}
\sisetup{locale=FR}
\usepackage[version=4]{mhchem}
\usepackage{listings}
\usepackage{fancyhdr}
\usepackage{lastpage}
\usepackage[most]{tcolorbox}
\usepackage{enumitem}
\setlist{nosep,leftmargin=*,topsep=2pt}
\usepackage[hidelinks,unicode]{hyperref}
\usepackage{bookmark}
\hypersetup{pdftitle={Banque PT — Langue vivante A 2026},pdfauthor={Rédigé par IA},pdfcreator={LaTeX},pdfsubject={Corrigé guidé}}
\usepackage{xurl}

% --- Pandoc : listes, tableaux, figures ----------------------------------------
\providecommand{\passthrough}[1]{#1}
\providecommand{\tightlist}{\setlength{\itemsep}{0pt}\setlength{\parskip}{0pt}}

\makeatletter
\def\maxwidth{\ifdim\Gin@nat@width>\linewidth\linewidth\else\Gin@nat@width\fi}
\def\maxheight{\ifdim\Gin@nat@height>0.55\textheight 0.55\textheight\else\Gin@nat@height\fi}
\makeatother
\setkeys{Gin}{width=\maxwidth,height=\maxheight,keepaspectratio}
\setcounter{secnumdepth}{-\maxdimen}
\providecommand{\pandocbounded}[1]{#1}
\usepackage{longtable,booktabs,array}

% --- Code (listings) -----------------------------------------------------------
\lstset{
  basicstyle=\ttfamily\small,
  columns=fullflexible,
  keepspaces=true,
  breaklines=true,
  breakatwhitespace=false,
  postbreak=\mbox{\textcolor{gris50}{$\hookrightarrow$}\space},
  frame=single,
  rulecolor=\color{gris35},
  framesep=4pt,
  xleftmargin=5pt,xrightmargin=5pt,
  backgroundcolor=\color{gris3},
  keywordstyle=\bfseries,
  commentstyle=\color{gris65},
  stringstyle=\color{gris80},
  showstringspaces=false,
  tabsize=4,
  extendedchars=true,
  aboveskip=6pt,belowskip=6pt,
}

% --- Macros des corrigés (KaTeX ↔ LaTeX) ---------------------------------------
% Les corrigés sont écrits pour KaTeX : on définit ici ce que KaTeX connaît et
% que LaTeX n'a pas, sans jamais retoucher le Markdown.
\providecommand{\R}{\mathbb{R}}
\providecommand{\N}{\mathbb{N}}
\providecommand{\Z}{\mathbb{Z}}
\providecommand{\Q}{\mathbb{Q}}
\providecommand{\C}{\mathbb{C}}
\providecommand{\dd}{\mathrm{d}}
\providecommand{\llbracket}{\lBrack}
\providecommand{\rrbracket}{\rBrack}
\providecommand{\lBrace}{\lbrace}
\providecommand{\rBrace}{\rbrace}
\providecommand{\Box}{\square}
\providecommand{\Diamond}{\lozenge}
\providecommand{\varnothing}{\emptyset}
\providecommand{\argmax}{\operatorname*{arg\,max}}
\providecommand{\argmin}{\operatorname*{arg\,min}}
\renewcommand{\leq}{\leqslant}
\renewcommand{\geq}{\geqslant}

% Équation hors texte : réduite à la largeur de la ligne si elle déborde.
\newsavebox{\boiteequation}
\newcommand{\equationajustee}[1]{%
  \sbox{\boiteequation}{$\displaystyle #1$}%
  \ifdim\wd\boiteequation>\linewidth
    \[\resizebox{\linewidth}{!}{\usebox{\boiteequation}}\]%
  \else
    \[#1\]%
  \fi}

% --- Encadrés pédagogiques (sobres, lisibles en noir et blanc) -----------------
\tcbset{
  encadre/.style={
    enhanced,breakable,sharp corners,
    colback=white,colframe=gris55,
    boxrule=0.4pt,leftrule=2.2pt,
    left=6pt,right=6pt,top=3pt,bottom=3pt,
    before skip=6pt,after skip=6pt,
    fonttitle=\bfseries\small,coltitle=black,
    attach title to upper,after title={\quad},
    parbox=false,
  }
}
\newtcolorbox{indication}[1]{encadre,colframe=gris45,colback=gris2,title={Indication~#1}}
\newtcolorbox{methode}{encadre,colframe=gris70,title={Méthode}}
\newtcolorbox{piege}{encadre,colframe=gris85,boxrule=0.8pt,leftrule=2.2pt,title={Piège}}
\newtcolorbox{jury}{encadre,colframe=gris60,colback=gris5,title={Jury}}
\newtcolorbox{encadrenote}{encadre,colframe=gris40}

% --- Titres de section ---------------------------------------------------------
\usepackage{titlesec}
\titleformat{\section}{\Large\bfseries}{}{0pt}{}
\titleformat{\subsection}{\large\bfseries}{}{0pt}{}
\titleformat{\subsubsection}{\normalsize\bfseries}{}{0pt}{}
\titlespacing*{\section}{0pt}{14pt plus 4pt minus 2pt}{6pt}
\titlespacing*{\subsection}{0pt}{10pt plus 3pt minus 2pt}{4pt}
\titlespacing*{\subsubsection}{0pt}{8pt plus 2pt minus 2pt}{2pt}

% --- En-tête et pied de page ---------------------------------------------------
\pagestyle{fancy}
\fancyhf{}
\fancyhead[L]{\small\textsc{Corrigé guidé} · Banque PT PT}
\fancyhead[R]{\small Langue vivante A 2026}
\fancyfoot[L]{\footnotesize Corrigé rédigé par IA, non officiel}
\fancyfoot[R]{\footnotesize \thepage\,/\,\pageref*{LastPage}}
\renewcommand{\headrulewidth}{0.4pt}
\renewcommand{\footrulewidth}{0pt}
\fancypagestyle{plain}{\fancyhf{}\fancyfoot[R]{\footnotesize \thepage\,/\,\pageref*{LastPage}}\fancyfoot[L]{\footnotesize Corrigé rédigé par IA, non officiel}\renewcommand{\headrulewidth}{0pt}}

\begin{document}
\thispagestyle{plain}
\begin{center}
  {\small\scshape Corrigé guidé}\par\smallskip
  {\LARGE\bfseries Banque PT — Langue vivante A 2026\par}\medskip
  {\large Banque PT · PT · Langue vivante A · 2026\par}
\end{center}
\vspace{2pt}\hrule\vspace{8pt}

\emph{Corrigé rédigé par IA, non officiel. Signalez toute erreur.}

\hypertarget{pruxe9sentation-du-sujet}{%
\section{Présentation du sujet}\label{pruxe9sentation-du-sujet}}

\textbf{Format.} L'épreuve dure 3 heures, sans dictionnaire ni appareil électronique. Le fascicule propose cinq langues (allemand, anglais, arabe, espagnol, italien) ; le candidat compose dans la seule langue choisie à l'inscription, faute de quoi il reçoit zéro. Chaque langue donne un dossier de cinq documents et une seule consigne :

\begin{longtable}[]{@{}
  >{\raggedright\arraybackslash}p{(\columnwidth - 2\tabcolsep) * \real{0.1916}}
  >{\raggedright\arraybackslash}p{(\columnwidth - 2\tabcolsep) * \real{0.7784}}@{}}
\toprule\noalign{}
\begin{minipage}[b]{\linewidth}\raggedright
Élément
\end{minipage} & \begin{minipage}[b]{\linewidth}\raggedright
Exigence
\end{minipage} \\
\midrule\noalign{}
\endhead
\bottomrule\noalign{}
\endlastfoot
Exercice & synthèse répondant à la question posée, « En vous appuyant uniquement sur les documents du dossier thématique qui vous est proposé » \\
Longueur & entre 450 et 500 mots \\
Titre & obligatoire, compté dans le nombre de mots \\
Décompte & nombre de mots (titre inclus) indiqué à la fin de la copie \\
\end{longtable}

Ce corrigé traite les cinq langues, soit 5 exercices numérotés de 1 à 5.

\textbf{Les cinq dossiers.}

\begin{longtable}[]{@{}
  >{\raggedright\arraybackslash}p{(\columnwidth - 6\tabcolsep) * \real{0.0720}}
  >{\raggedright\arraybackslash}p{(\columnwidth - 6\tabcolsep) * \real{0.2668}}
  >{\raggedright\arraybackslash}p{(\columnwidth - 6\tabcolsep) * \real{0.3356}}
  >{\raggedright\arraybackslash}p{(\columnwidth - 6\tabcolsep) * \real{0.2956}}@{}}
\toprule\noalign{}
\begin{minipage}[b]{\linewidth}\raggedright
Langue
\end{minipage} & \begin{minipage}[b]{\linewidth}\raggedright
Question
\end{minipage} & \begin{minipage}[b]{\linewidth}\raggedright
Textes
\end{minipage} & \begin{minipage}[b]{\linewidth}\raggedright
Documents non textuels
\end{minipage} \\
\midrule\noalign{}
\endhead
\bottomrule\noalign{}
\endlastfoot
Allemand & \emph{Ist eine Welt ohne Diskriminierungen immer noch ein Traum?} & loi générale sur l'égalité de traitement (AGG) ; antitsiganisme à l'école (\emph{Tagesschau}, 2025) ; discriminations produites par l'intelligence artificielle (\emph{Sonntagsblatt}, 2024) & infographie d'un cabinet d'avocats sur le harcèlement sexuel au travail ; photo de dés, l'un portant un pictogramme de fauteuil roulant \\
Anglais & \emph{To what extent can geoengineering contribute to the fight against climate change?} & \emph{The Economist} (2023) ; \emph{The Guardian} (2024) ; \emph{BBC News} (2024) & schéma des techniques de géo-ingénierie (CFR Education, 2024) ; dessin satirique « Problem Inc.~/ Solution Inc.~» (2016) \\
Arabe & Dans quelle mesure le développement durable est-il une priorité pour l'avenir des populations du monde arabe ? & \emph{Jadaliyya} (2024) ; \emph{Al Jazeera} (2019) ; \emph{BBC Arabic} (2023) & infographie de \emph{Sky News Arabia} sur l'usage de l'eau par pays (2015) ; affiche des 17 objectifs de développement durable (2022) \\
Espagnol & \emph{¿De qué manera está evolucionando el modelo turístico español?} & \emph{El País}, \emph{eldiario.es}, \emph{zonaestrategia.net} (2024) & deux dessins : « ¿Turismofóbicos? » (Padylla) et « 13, rue de la Turistificación » \\
Italien & \emph{In che misura l'Italia subisce e fa fronte al sovraffollamento turistico?} & \emph{leomagazineofficial.it}, \emph{ansa.it}, \emph{italiaatavola.net} (2024) & tableau Istat des nuitées européennes par région (2022) ; photo de la foule devant la fontaine de Trevi \\
\end{longtable}

\textbf{Ce qui est évalué.} Le rapport 2026 ne publie un barème détaillé qu'en anglais : « langue, méthodologie, compréhension/ restitution, respectivement sur 20, 15 et 15 points (total sur 50) » (p.~4). Les cinq jurys convergent sur les mêmes exigences :

\begin{itemize}
\tightlist
\item
  un titre (son oubli entraîne un malus partout) et le nombre de mots en fin de copie ;
\item
  aucune opinion personnelle, aucune connaissance extérieure au dossier, y compris dans l'accroche ;
\item
  chaque document cité explicitement (« doc. 1 », « document 2 »\ldots), images et tableaux compris, et réellement analysé ;
\item
  une reformulation des idées, sans recopier les textes ;
\item
  une introduction courte (situer le thème, poser la problématique, annoncer le plan), un développement en deux ou trois parties équilibrées, une conclusion soignée.
\end{itemize}

\textbf{Difficulté.} Les dossiers 2026 étaient, de l'aveu des jurys, accessibles. La sélection s'est donc faite sur la précision de la restitution (contresens sur un chiffre, une nuance, une ironie), sur le traitement des documents iconographiques, souvent expédiés, et sur la qualité de la langue.

\textbf{Gestion du temps.} Une répartition raisonnable : 40 minutes de lecture et de grille d'analyse (une ligne par document : idée principale, chiffres, ton), 20 minutes de plan, 1 h 40 de rédaction directe, 20 minutes de relecture (accords, conjugaisons) et de décompte. Le jury d'allemand conseille de marquer un repère tous les 50 mots pour faciliter le contrôle.

\textbf{Statistiques publiées.} Anglais : 2610 copies (p.~4), sans moyenne. Arabe : 22 candidats, notes de 10,80 à 18,80, moyenne 14,47 (p.~10). Espagnol : 23 candidats, moyenne 11,15, notes de 5,6 à 18,4 (p.~13). Italien : cinq candidats (p.~16). Allemand : pas de chiffre, mais « peu de notes inférieures à 10 cette année » (p.~1).

\textbf{Règle de décompte.} Le sujet ne la précise pas. Ce corrigé compte les suites de caractères séparées par des espaces, titre inclus ; un nombre (\emph{8,6}, \emph{1.2°C}) compte pour un mot, un signe isolé (\%) pour zéro. En italien, une forme élidée (\emph{l'Italia}) compte pour un mot ; même en la comptant pour deux, la synthèse italienne reste sous 500 mots. En arabe, on compte les mots graphiques (و، ب، ل، ال soudés ne sont pas des mots à part).

\hypertarget{partie-i-allemand}{%
\section{Partie I --- Allemand}\label{partie-i-allemand}}

Dossier sur les discriminations : leurs formes (au travail, à l'école, dans les algorithmes, envers les personnes handicapées) et les moyens d'y remédier (loi, éducation, technique, démarche individuelle).

\hypertarget{question-1-allemand-synthuxe8se}{%
\subsection{Question 1 --- Allemand, synthèse}\label{question-1-allemand-synthuxe8se}}

\begin{indication}{1}

Avant tout plan, dressez une grille : pour chaque document, l'idée principale et les chiffres. Classez ensuite chaque idée en « constat » ou en « remède ». Les documents 4 et 5 ont autant de poids que les textes.

\end{indication}

\begin{indication}{2}

\emph{immer noch ein Traum?} appelle une réponse nuancée : la discrimination recule-t-elle, ou change-t-elle de forme ? Cherchez ce que les documents 1, 2 et 3 ont en commun (une discrimination souvent involontaire).

\end{indication}

\begin{indication}{3}

La photo du document 5 ne se décrit pas seulement : interprétez le geste de la main qui tient le dé au-dessus de la pyramide. Le mot allemand pour « dé » est un nom masculin fréquent.

\end{indication}

\hypertarget{solution}{%
\subsubsection{Solution}\label{solution}}

\textbf{Analyse du dossier.}

\begin{longtable}[]{@{}
  >{\raggedright\arraybackslash}p{(\columnwidth - 6\tabcolsep) * \real{0.0769}}
  >{\raggedright\arraybackslash}p{(\columnwidth - 6\tabcolsep) * \real{0.2339}}
  >{\raggedright\arraybackslash}p{(\columnwidth - 6\tabcolsep) * \real{0.4631}}
  >{\raggedright\arraybackslash}p{(\columnwidth - 6\tabcolsep) * \real{0.1961}}@{}}
\toprule\noalign{}
\begin{minipage}[b]{\linewidth}\raggedright
Doc.
\end{minipage} & \begin{minipage}[b]{\linewidth}\raggedright
Nature
\end{minipage} & \begin{minipage}[b]{\linewidth}\raggedright
Idée principale
\end{minipage} & \begin{minipage}[b]{\linewidth}\raggedright
Place dans le plan
\end{minipage} \\
\midrule\noalign{}
\endhead
\bottomrule\noalign{}
\endlastfoot
1 & texte (rightmart.de) & discrimination au travail (un salarié sur trois) ; l'AGG de 2006 étend l'interdiction aux employeurs privés ; discrimination directe ou indirecte & I (ampleur) et II (droit) \\
2 & texte (\emph{Tagesschau}, 2025) & antitsiganisme institutionnel à l'école, préjugés involontaires des enseignants ; conseillers d'éducation, enseignement culturellement adapté & I et II (éducation) \\
3 & texte (\emph{Sonntagsblatt}, 2024) & l'IA reproduit et amplifie les biais de ses données ; agir dès la programmation, responsabilité des États et des entreprises & I et II (technique) \\
4 & infographie (cabinet d'avocats, 2023) & conduite à tenir face au harcèlement sexuel : évaluer, parler, informer l'employeur, consulter & II (démarche individuelle) \\
5 & photo (\emph{Ärzteblatt}) & personnes handicapées et malades chroniques : leur place dans la société dépend des autres & I (publics touchés) \\
\end{longtable}

\textbf{Plan.} I. Une discrimination omniprésente et souvent invisible (docs 1, 2, 3, 5). II. Des remèdes juridiques, éducatifs et techniques (docs 1, 4, 2, 3). Conclusion : le rêve n'est pas atteint, mais il n'est plus hors de portée.

\textbf{Gleich vor dem Gesetz, ungleich im Alltag}

Das Dossier besteht aus drei Presseartikeln, einer Infografik und einem Foto, die sich mit Diskriminierung im Berufsleben, in der Schule und in der digitalen Welt beschäftigen. Ist eine Gesellschaft ohne Benachteiligung erreichbar oder bleibt sie eine Utopie? Zunächst wird gezeigt, wie verbreitet Diskriminierung heute ist, danach, welche Gegenmittel es gibt.

Erstens geht aus dem Dossier hervor, dass Diskriminierung keineswegs der Vergangenheit angehört. Laut Dokument 1 gibt jeder dritte Arbeitnehmer in Deutschland an, in jüngerer Zeit selbst benachteiligt worden zu sein, vor allem wegen seines Alters oder seines Geschlechts. Fast die Hälfte der Fälle ereignet sich im Arbeitsleben. Besonders gefährlich ist nach Dokument 2 die unbewusste, mittelbare Diskriminierung: Sinti und Roma erleben an Schulen nicht nur Beleidigungen, sondern einen institutionellen Antiziganismus. Ihre Kinder wurden oft pauschal auf Förderschulen geschickt, und eine Studie zeigt, dass angehende Lehrkräfte Kindern mit Migrationshintergrund bei gleichen Leistungen seltener das Gymnasium empfehlen. Die 484 seit 2023 gemeldeten Vorfälle im Bildungsbereich sind dabei nur die Spitze des Eisbergs. Dokument 3 ergänzt, dass selbst angeblich objektive Maschinen nicht neutral sind: Da künstliche Intelligenz aus verzerrten Daten lernt, erkennt sie dunkle Gesichter schlechter, bevorzugt Männer bei der Personalauswahl und gewährt Frauen niedrigere Kredite. Sie kann bestehende Vorurteile also sogar verstärken. Schließlich erinnert das Foto in Dokument 5 daran, dass auch Menschen mit Behinderung oder chronischer Krankheit betroffen sind. Eine Hand hält einen Würfel mit einem Rollstuhlsymbol über einer Pyramide aus Würfeln, die Männer, Frauen und Familien zeigen: Ob dieser Baustein seinen Platz in der Gesellschaft findet, hängt von den anderen ab.

Zweitens zeigen die Dokumente jedoch, dass es wirksame Gegenmittel gibt. Das wichtigste ist das Recht. Während Artikel 3 des Grundgesetzes vor allem den Staat bindet, verbietet das Allgemeine Gleichbehandlungsgesetz seit 2006 auch privaten Arbeitgebern jede Benachteiligung wegen der Herkunft, des Geschlechts, der Religion, einer Behinderung, des Alters oder der sexuellen Identität (Dokument 1). Betroffene können Ansprüche gegen Arbeitgeber, Vorgesetzte und sogar Kollegen geltend machen. Wie man konkret vorgeht, veranschaulicht die Infografik in Dokument 4 am Beispiel der sexuellen Belästigung: Das Opfer soll zunächst entscheiden, was richtig und was falsch ist, gegebenenfalls das Gespräch suchen, dann den Arbeitgeber informieren und sich notfalls an eine Kanzlei wenden. Recht allein genügt aber nicht, denn tief verwurzelte Vorurteile lassen sich nicht verbieten. Deshalb setzt Dokument 2 auf Bildung: Bildungsberater sollen eine gleichberechtigte Teilhabe sichern, und junge Lehrkräfte sollen lernen, die kulturelle Herkunft der Kinder in den Unterricht einzubeziehen. Für die künstliche Intelligenz fordert Dokument 3, schon bei der Programmierung anzusetzen und nur korrekte Daten zu verwenden. Letztlich tragen aber Regierungen und Technologiekonzerne die Verantwortung, während jeder Nutzer die Ergebnisse kritisch prüfen sollte.

Zusammenfassend lässt sich sagen, dass eine Welt ohne Diskriminierung noch nicht Wirklichkeit ist. Die Benachteiligung hat lediglich ihre Formen verändert und versteckt sich heute in Institutionen und Algorithmen. Die Dokumente zeigen aber auch, dass Gesetze, Bildung und Wachsamkeit diesen Traum Schritt für Schritt näher rücken lassen.

\textbf{(482 Wörter, Titel inbegriffen)}

Choix justifiés :

\begin{itemize}
\tightlist
\item
  Le titre oppose deux adjectifs symétriques (\emph{gleich / ungleich}) et résume la tension du dossier sans recopier la question.
\item
  \emph{Laut Dokument 1 gibt jeder dritte Arbeitnehmer an, \ldots{} benachteiligt worden zu sein} : infinitif passé passif, qui reformule le texte au lieu de le citer.
\item
  \emph{Während} oppose deux propositions (verbe en fin de subordonnée), \emph{Da} introduit une cause connue, \emph{Ob \ldots{} findet} une interrogative indirecte : trois subordonnées différentes, valorisées par le jury.
\item
  La conclusion répond à la question (\emph{noch nicht Wirklichkeit}) sans opinion personnelle : c'est le dossier qui « montre » (\emph{Die Dokumente zeigen}).
\end{itemize}

\begin{methode}

Pour un document iconographique : décrire en une phrase (\emph{Eine Hand hält einen Würfel \ldots{}}), puis interpréter en une phrase ce que l'image signifie au regard de la question. La description seule ne rapporte presque rien.

\end{methode}

\begin{piege}

Genres et formes à ne pas manquer : \emph{das Dossier}, \emph{die Diskriminierung}, \emph{die Benachteiligung}, \emph{das Vorurteil}, \emph{der Würfel}, \emph{das Gesetz}. \emph{benachteiligen} se construit avec l'accusatif (\emph{Frauen werden benachteiligt}). \emph{Ich finde}, \emph{meiner Meinung nach} sont exclus d'une synthèse.

\end{piege}

\begin{jury}

Rapport 2026 (p.~1) : les documents étaient à analyser « de manière équivalente en termes de nombre de mots », et il est « strictement interdit d'y introduire un avis personnel, un jugement de valeur, ou toute information extérieure au corpus ». Le jury note que « très peu connaissaient le mot allemand pour « dé », utile pour décrire le document 5 », et rappelle qu'« Une introduction qui fait 15 lignes et une conclusion qui n'en fait qu'une déséquilibre le travail. » Côté forme (p.~2) : « Un décompte tous les 50 mots est une bonne option pour aider au contrôle. » et, pour l'écriture illisible, « Dans le doute, le jury tranchera en faveur d'une faute. »

\end{jury}

\hypertarget{partie-ii-anglais}{%
\section{Partie II --- Anglais}\label{partie-ii-anglais}}

Dossier technique sur la géo-ingénierie : le problème du réchauffement, la palette des techniques (renvoyer le rayonnement solaire, capter le carbone), leurs risques, et l'appel à la recherche.

\hypertarget{question-2-anglais-synthuxe8se}{%
\subsection{Question 2 --- Anglais, synthèse}\label{question-2-anglais-synthuxe8se}}

\begin{indication}{1}

Le schéma du document 4 donne la carte de toutes les techniques : partez de lui pour classer les informations dispersées dans les textes (réflexion du rayonnement solaire d'un côté, absorption du carbone de l'autre).

\end{indication}

\begin{indication}{2}

Relisez avec soin trois passages : la comparaison des coûts (document 1), la position exacte de Richard Spinrad et le modal \emph{may} (document 2), le chiffre de la capture directe rapporté aux émissions mondiales (document 3).

\end{indication}

\begin{indication}{3}

Le dessin de 2016 est ironique : que suggère le tuyau qui relie « Problem Inc.~» à « Solution Inc.~» ? Et le dossier ne se clôt pas sur un oui ou un non : quel troisième axe traverse les documents 2 et 3 ?

\end{indication}

\hypertarget{solution-1}{%
\subsubsection{Solution}\label{solution-1}}

\textbf{Analyse du dossier.}

\begin{longtable}[]{@{}
  >{\raggedright\arraybackslash}p{(\columnwidth - 6\tabcolsep) * \real{0.0724}}
  >{\raggedright\arraybackslash}p{(\columnwidth - 6\tabcolsep) * \real{0.1880}}
  >{\raggedright\arraybackslash}p{(\columnwidth - 6\tabcolsep) * \real{0.5439}}
  >{\raggedright\arraybackslash}p{(\columnwidth - 6\tabcolsep) * \real{0.1658}}@{}}
\toprule\noalign{}
\begin{minipage}[b]{\linewidth}\raggedright
Doc.
\end{minipage} & \begin{minipage}[b]{\linewidth}\raggedright
Nature
\end{minipage} & \begin{minipage}[b]{\linewidth}\raggedright
Idée principale
\end{minipage} & \begin{minipage}[b]{\linewidth}\raggedright
Place dans le plan
\end{minipage} \\
\midrule\noalign{}
\endhead
\bottomrule\noalign{}
\endlastfoot
1 & \emph{The Economist}, 2023 & la géo-ingénierie devient respectable ; mécanisme de l'effet de serre (280 → 417 ppm, +1,2 °C) ; Pinatubo ; quelques dizaines de milliards par an contre des milliers de milliards pour décarboner ; effets secondaires incertains & I, II \\
2 & \emph{The Guardian}, 2024 & la NOAA étudie les effets ; Spinrad : certaines techniques seront probablement utiles, mais il faut d'abord connaître leurs conséquences ; financement menacé ; les Américains, qui paient les catastrophes, \emph{may be more open} & I (conséquences), III \\
3 & \emph{BBC News}, 2024 & industrie en croissance et théories du complot ; capture directe : 10 000 t par an contre 35 milliards émises ; éclaircissement des nuages : sécheresse possible au Brésil ; risque de détourner de la réduction des émissions & II, III \\
4 & schéma, CFR Education, 2024 & panorama : nuages marins, aérosols stratosphériques, miroirs spatiaux ; fertilisation des océans, minéralisation, capture directe, reboisement & II \\
5 & dessin, 2016 & le CO₂ de « Problem Inc.~» et le SO₂ de « Solution Inc.~» relèvent de la même logique : le remède semble aussi nocif que le mal & III \\
\end{longtable}

\textbf{Plan.} I. Un réchauffement d'origine humaine aux effets déjà coûteux. II. Une palette de techniques prometteuses et bon marché. III. Des risques mal connus, d'où l'appel à la recherche. Conclusion : un complément possible, pas un substitut à la réduction des émissions.

\textbf{Cooling the Planet: Shortcut or Short Circuit?}

Since the Industrial Revolution, human activity has steadily warmed the planet. Geoengineering, the deliberate manipulation of the climate to cool the Earth, is increasingly presented as a possible response. Drawing on three press articles, an infographic and a cartoon, this synthesis examines what these techniques could achieve, what risks they entail and why research is needed.

According to document 1, greenhouse gases such as carbon dioxide trap part of the heat that the Earth re-emits. Their concentration has risen from about 280 to 417 parts per million since pre-industrial times, raising average temperatures by roughly 1.2°C. The consequences are already tangible: according to document 2, Americans have suffered wildfires and a record number of costly disasters, which they now pay for through their insurance.

Geoengineering covers a wide range of responses, summed up in the infographic of document 4. Some techniques reflect sunlight back into space: stratospheric aerosols, marine cloud brightening, which consists in spraying fine seawater towards low clouds (document 3), or even space mirrors. Others absorb carbon dioxide, through ocean fertilization, carbon mineralization, direct air capture or reforestation. Solar geoengineering imitates volcanoes: Mount Pinatubo cooled the planet by up to 0.5°C in 1991. Its main appeal is its cost. Injecting about two million tonnes of sulphur a year to lower temperatures by 1°C might cost a few tens of billions of dollars annually, whereas decarbonising the world economy would cost trillions (document 1).

Yet side effects remain largely unknown: according to modelling quoted in document 3, brightening clouds off Namibia could cause droughts in Brazil and threaten its rainforests, while some regions might even get warmer. Direct air capture, the most advanced technique, removes only about 10,000 tonnes of carbon dioxide a year, compared with some 35 billion tonnes emitted worldwide. Experts also fear that such a quick fix may distract governments from reducing emissions. The 2016 cartoon in document 5 takes this criticism further: a factory labelled ``Problem Inc.'' releases carbon dioxide, while ``Solution Inc.'', connected to it by a pipe, emits sulphur dioxide. The irony is clear: the remedy appears as harmful as the disease.

This is why the documents call for research rather than hasty deployment. Richard Spinrad, head of the US agency NOAA, suspects that some techniques will be part of the solution, but insists that their consequences for ecosystems must first be understood (document 2). He also defends his agency's funding, which some members of Congress threaten to cut, and notes that Americans, increasingly aware of the cost of climate change, may be more open to new technologies. Rigorous science could also counter the conspiracy theories that BBC Weather has seen spreading online (document 3).

Geoengineering can therefore contribute to the fight against climate change, but only as a complement to emission cuts. Cheap and increasingly respectable, it will remain a risky gamble as long as its side effects have not been properly assessed.

\textbf{(488 words, title included)}

Choix justifiés :

\begin{itemize}
\tightlist
\item
  Le titre joue sur \emph{shortcut} (raccourci) et \emph{short circuit} (court-circuit) : il pose l'alternative du dossier sans la formule toute faite \emph{a boon or a bane}.
\item
  \emph{Since the Industrial Revolution, human activity has steadily warmed the planet} : \emph{present perfect} avec \emph{since}, l'erreur que le jury a relevée dans la majorité des introductions.
\item
  \emph{space mirrors} est nommé : c'est la seule façon de prouver que le document 4 a été lu.
\item
  Le coût est restitué dans le bon sens (\emph{a few tens of billions \ldots{} whereas \ldots{} trillions}) et Spinrad reste nuancé (\emph{suspects}, \emph{must first be understood}, \emph{may be more open}).
\item
  \emph{global warming} et \emph{climate change} sans article ; \emph{research} indénombrable.
\end{itemize}

\begin{methode}

Une accroche d'introduction se construit avec une idée du dossier reformulée (ici la hausse des températures du document 1), jamais avec une citation entre guillemets ni une connaissance extérieure.

\end{methode}

\begin{piege}

Fautes relevées par le jury : \emph{to what extent can geoengineering contribute} dans une interrogative indirecte (il faut \emph{to what extent geoengineering can contribute}), \emph{researches}, \emph{damages}, \emph{environmentals problems}, \emph{faced to} (au lieu de \emph{faced with}), \emph{the global warming}, et \emph{important} au sens quantitatif (\emph{a significant amount}, \emph{a large number}).

\end{piege}

\begin{jury}

Rapport 2026 : il était difficile de savoir si le document 4 était traité « à moins de mentionner « space mirrors » » (p.~3). Sur le coût, « Beaucoup de candidats ont en effet restitué l'inverse : ces technologies seraient coûteuses. » ; sur Spinrad, « sa position était beaucoup plus nuancée » ; sur le dessin, « Le jury attendait a minima un commentaire sur le fait que la solution semblait aussi nocive que le problème. » (p.~4). Idée peu restituée et bonifiée : « la nécessité de la recherche » (p.~5). Enfin, « Le jury attend une introduction courte, permettant de situer les enjeux du dossier. » (p.~5) et « Le titre est le seul espace de créativité et de liberté dans cette épreuve de synthèse » (p.~6).

\end{jury}

\hypertarget{partie-iii-arabe}{%
\section{Partie III --- Arabe}\label{partie-iii-arabe}}

Dossier sur le développement durable dans le monde arabe : sa définition et ses conditions politiques, la crise de l'eau comme urgence principale, la faiblesse du journalisme environnemental.

\hypertarget{question-3-arabe-synthuxe8se}{%
\subsection{Question 3 --- Arabe, synthèse}\label{question-3-arabe-synthuxe8se}}

\begin{indication}{1}

La question porte sur une « priorité » : cherchez dans le dossier ce qui rend le développement durable urgent (pauvreté, eau) et ce qui conditionne sa réalisation (institutions, information).

\end{indication}

\begin{indication}{2}

Les documents 3 et 4 parlent tous deux de l'eau : regroupez leurs chiffres dans une même partie. Les documents 1 et 5 se répondent aussi (objectifs du développement durable).

\end{indication}

\begin{indication}{3}

Pour l'infographie du document 4, retenez deux catégories de pays et le sens général (consommation supérieure aux ressources), sans recopier tous les pourcentages.

\end{indication}

\hypertarget{solution-2}{%
\subsubsection{Solution}\label{solution-2}}

\textbf{Analyse du dossier.}

\begin{longtable}[]{@{}
  >{\raggedright\arraybackslash}p{(\columnwidth - 6\tabcolsep) * \real{0.0708}}
  >{\raggedright\arraybackslash}p{(\columnwidth - 6\tabcolsep) * \real{0.2095}}
  >{\raggedright\arraybackslash}p{(\columnwidth - 6\tabcolsep) * \real{0.5394}}
  >{\raggedright\arraybackslash}p{(\columnwidth - 6\tabcolsep) * \real{0.1502}}@{}}
\toprule\noalign{}
\begin{minipage}[b]{\linewidth}\raggedright
Doc.
\end{minipage} & \begin{minipage}[b]{\linewidth}\raggedright
Nature
\end{minipage} & \begin{minipage}[b]{\linewidth}\raggedright
Idée principale
\end{minipage} & \begin{minipage}[b]{\linewidth}\raggedright
Place dans le plan
\end{minipage} \\
\midrule\noalign{}
\endhead
\bottomrule\noalign{}
\endlastfoot
1 & texte (\emph{Jadaliyya}, 2024) & développement durable = économie, social, environnement ; impossible sans justice, institutions responsables, libertés ; pauvreté (environ 30 \% des Égyptiens), eau, assainissement, énergie solaire & I \\
2 & texte (\emph{Al Jazeera}, 2019) & les médias arabes couvrent les catastrophes sans en analyser causes et responsabilités ; retard, traductions, absence de financement ; triple regard social, environnemental, économique & III \\
3 & texte (\emph{BBC Arabic}, 2023) & 390 millions d'Arabes (90 \%) en pénurie d'eau ; Irak, Égypte (500 m³ par habitant), Maroc (2500 → 650 m³) ; coupures en Tunisie, dessalement en Algérie, récupération des pluies en Arabie saoudite ; changer la culture de consommation & II \\
4 & infographie (\emph{Sky News Arabia}, 2015) & région parmi les plus pauvres en eau ; usage supérieur aux ressources, surtout dans les pays pétroliers (Émirats, Arabie saoudite) & II \\
5 & affiche (2022) & 17 objectifs interdépendants à atteindre d'ici 2030 pour ne laisser personne de côté & I \\
\end{longtable}

\textbf{Plan.} I. Un développement global qui exige justice et démocratie (docs 1, 5). II. L'eau, priorité la plus urgente (docs 3, 4). III. Un enjeu qui suppose une information de qualité (doc. 2). Conclusion : une nécessité vitale, à des conditions politiques, techniques et culturelles.

\textbf{التنمية المستدامة في العالم العربي: ضرورة لا ترف}

يضم هذا الملف خمس وثائق متنوعة: ثلاثة مقالات صحفية نشرتها مواقع جدلية والجزيرة وبي بي سي عربي، ورسم بياني من موقع سكاي نيوز عربية حول استخدام المياه، وملصق يعرض أهداف التنمية المستدامة. وتطرح هذه الوثائق سؤالا جوهريا: إلى أي مدى تشكل التنمية المستدامة أولوية لمستقبل الشعوب العربية؟ سنعرض أولا مفهوم التنمية المستدامة وشروطها، ثم نتوقف عند أزمة المياه بوصفها أخطر التحديات، وأخيرا عند دور الإعلام في مواجهتها.

تؤكد الوثيقة الأولى أن التنمية المستدامة مفهوم شامل يجمع بين النمو الاقتصادي والتقدم الاجتماعي وحماية البيئة. وترى الكاتبة أن تحقيقها مستحيل في ظل الفساد والفقر وغياب الحريات، لأنها تحتاج إلى مؤسسات فعالة خاضعة للمساءلة، وإلى العدل والمساواة وسيادة القانون، ولذلك تربط بينها وبين الديمقراطية. ومن أولويات هذه التنمية القضاء على الفقر، إذ يعيش نحو 30 في المائة من المصريين تحت خط الفقر، إلى جانب ضمان الماء والصرف الصحي والطاقة النظيفة بأسعار معقولة، كألواح الطاقة الشمسية في القرى النائية. وتظهر هذه الأهداف نفسها في الوثيقة الخامسة التي تعرض أهداف التنمية المستدامة السبعة عشر، من القضاء على الفقر والجوع إلى العمل المناخي والسلام والعدل، وتشدد على ترابطها وعلى ضرورة تحقيقها جميعا بحلول عام 2030 حتى لا يتخلف أحد عن الركب.

غير أن الماء يبدو التحدي الأكثر إلحاحا. فحسب الوثيقة الثالثة، يعيش 390 مليون عربي، أي نحو 90 في المائة من السكان، في دول تعاني من ندرة المياه، وتقع 19 دولة عربية من أصل 22 في نطاق الشح المائي. فالعراق، بلاد الرافدين، صار يعاني من الجفاف، ولا يتجاوز نصيب الفرد في مصر 500 متر مكعب سنويا، أي نصف حد الفقر المائي، كما تراجع نصيب الفرد في المغرب من 2500 متر مكعب عام 1960 إلى نحو 650 اليوم. وتؤكد الوثيقة الرابعة خطورة الوضع، فالمنطقة العربية من أفقر مناطق العالم مائيا، وكثير من دولها يستخدم من المياه أكثر مما تتيحه مواردها، ولا سيما الدول المصدرة للنفط كالإمارات والسعودية، وكذلك الأردن ومصر واليمن. وتحاول الحكومات الاستجابة لهذه الأزمة: فقد قطعت تونس الماء سبع ساعات يوميا، ودعا الرئيس الجزائري إلى إنشاء محطات لتحلية مياه البحر، وأقامت السعودية مشاريع لحصاد مياه الأمطار. لكن الخبراء يرون أن ذلك لا يكفي من دون تغيير جذري في ثقافة الاستهلاك والتخلي عن الإسراف.

ولا يمكن مواجهة هذه التحديات من دون إعلام قادر على توعية الجمهور. إلا أن الوثيقة الثانية تكشف ضعف الصحافة البيئية العربية، فهي تهتم بالكوارث كالسيول وانجراف التربة وبخسائرها المادية والبشرية، لكنها تهمل أسبابها والجهات المسؤولة عنها وآثارها على مستقبل السكان، بسبب نقص المعلومات. ويرى الباحث عمر عاصي أن هذا الإعلام ما زال متأخرا جدا، لأنه يعتمد على ترجمة مصادر أجنبية ويكتفي بوجهة نظر واحدة، فضلا عن غياب التمويل والاهتمام الرسمي. ويدعو إلى تناول كل قضية من زوايا ثلاث، اجتماعية وبيئية واقتصادية، إذ لا يجوز مثلا الإشادة بمشروع للطاقة الشمسية في بلد يستغل عمل الأطفال.

وخلاصة القول إن الوثائق تجمع على أن التنمية المستدامة ليست ترفا بل أولوية حيوية لمستقبل العالم العربي، ولا سيما في ظل الفقر وندرة المياه. غير أن تحقيقها يتطلب مؤسسات عادلة وديمقراطية، وسياسات مائية رشيدة، وإعلاما واعيا، ومواطنين مسؤولين.

\textbf{(478 كلمة، بما فيها العنوان)}

Choix justifiés :

\begin{itemize}
\tightlist
\item
  Titre court qui répond à la question (ضرورة لا ترف) sans la recopier.
\item
  Problématique posée dès l'introduction (إلى أي مدى تشكل التنمية المستدامة أولوية\ldots؟) et plan annoncé (أولا، ثم، وأخيرا).
\item
  Connecteurs variés : غير أن، فحسب، إلا أن، فضلا عن، وخلاصة القول.
\item
  Hamza de liaison sans hamza écrite (الاستجابة، الاستهلاك، اجتماعية، استخدام), إنّ après القول (وخلاصة القول إن), للنفط après المصدرة.
\end{itemize}

\begin{methode}

Les chiffres se regroupent par idée : trois ou quatre données bien choisies (90 \%, 500 m³, 2500 → 650 m³) valent mieux qu'une liste. Dans l'infographie, on retient la tendance et un ou deux exemples.

\end{methode}

\begin{piege}

Erreurs de langue relevées par le jury : hamza écrite à tort en début de mot (\emph{إستهلاك}, \emph{إعتبار}), orthographe de mots fréquents (ذلك، لكن، الجزائر، المغرب، التي), prépositions mal choisies (الدول المصدرة للنفط et non بالنفط ; في حاجة إلى et non من), accords (توجد مشاكل كثيرة ; أهداف التنمية المستدامة هي), calque de « où » relatif (حيث et non أين).

\end{piege}

\begin{jury}

Rapport 2026 : « Deux candidats n'ont pas attribué de titre à leur production, ce qui a entraîné une pénalité de points. » (p.~10). « Il est en effet essentiel de construire la synthèse autour d'une problématique bien formulée. » (p.~11). La conclusion du rapport demande aux candidats de veiller « à restituer fidèlement le contenu du dossier » et d'accorder « une attention particulière à la qualité de la langue arabe » (p.~12).

\end{jury}

\hypertarget{partie-iv-espagnol}{%
\section{Partie IV --- Espagnol}\label{partie-iv-espagnol}}

Dossier sur le tourisme de masse en Espagne : un secteur record, mais contesté pour ses effets sur le logement, l'emploi et l'environnement, et l'émergence de propositions de décroissance.

\hypertarget{question-4-espagnol-synthuxe8se}{%
\subsection{Question 4 --- Espagnol, synthèse}\label{question-4-espagnol-synthuxe8se}}

\begin{indication}{1}

La question porte sur une évolution : repérez ce qui change (chiffres records, nouveau cycle de protestations, propositions de limites) par rapport à ce qui dure (poids économique).

\end{indication}

\begin{indication}{2}

Les deux dessins sont humoristiques : expliquez le procédé (jeu de mots \emph{turismofóbicos / canariasfílicos} ; parodie de \emph{13, Rue del Percebe} où chaque appartement est devenu une location touristique) et ce qu'il dénonce.

\end{indication}

\begin{indication}{3}

Même le lobby du secteur (document 2) admet des limites : utilisez ce point pour nuancer et relier les documents entre eux.

\end{indication}

\hypertarget{solution-3}{%
\subsubsection{Solution}\label{solution-3}}

\textbf{Analyse du dossier.}

\begin{longtable}[]{@{}
  >{\raggedright\arraybackslash}p{(\columnwidth - 6\tabcolsep) * \real{0.0755}}
  >{\raggedright\arraybackslash}p{(\columnwidth - 6\tabcolsep) * \real{0.2200}}
  >{\raggedright\arraybackslash}p{(\columnwidth - 6\tabcolsep) * \real{0.5145}}
  >{\raggedright\arraybackslash}p{(\columnwidth - 6\tabcolsep) * \real{0.1601}}@{}}
\toprule\noalign{}
\begin{minipage}[b]{\linewidth}\raggedright
Doc.
\end{minipage} & \begin{minipage}[b]{\linewidth}\raggedright
Nature
\end{minipage} & \begin{minipage}[b]{\linewidth}\raggedright
Idée principale
\end{minipage} & \begin{minipage}[b]{\linewidth}\raggedright
Place dans le plan
\end{minipage} \\
\midrule\noalign{}
\endhead
\bottomrule\noalign{}
\endlastfoot
1 & tribune (\emph{El País}, 2024) & une industrie hors de contrôle expulse les habitants ; Cadix : chômage au-dessus de 20 \%, La Viña a perdu 20 \% de sa population ; eau, posidonie ; « les villes sont faites pour y vivre » & II, III \\
2 & article (\emph{eldiario.es}, 2024) & record annoncé (+8,6 \%, plus de 200 milliards d'euros) ; Exceltur reconnaît des limites et accuse les logements touristiques & I, II \\
3 & article (\emph{zonaestrategia.net}, 2024) & nouveau cycle de protestation (Canaries, 20 avril 2024) ; racines dans les crises de 2008 et du Covid ; décroître et transformer, risque d'élitisation & I, III \\
4 & dessin (Padylla, 2024) & les manifestants ne sont pas « turismofóbicos » mais « canariasfílicos » ; pancartes : tourisme durable, moratoire, écotaxe & III \\
5 & dessin (nuevatribuna.es, 2024) & immeuble entièrement converti (Airbnb, Booking, Starbucks) ; le résident relégué dans une mansarde de 20 m² à 700 € & II \\
\end{longtable}

\textbf{Plan.} I. Un secteur économique fondamental. II. Un modèle de plus en plus contesté pour ses effets sociaux, urbains et environnementaux. III. La nécessité de repenser le tourisme. C'est le plan suggéré par le jury.

\textbf{Turismo en España: del récord al límite}

El dossier reúne tres artículos de prensa de 2024 y dos viñetas satíricas que describen un turismo de récord frente al creciente malestar de los vecinos. Cabe preguntarse, por lo tanto, en qué medida el modelo turístico español está pasando de la expansión sin freno a la búsqueda de límites. Se analizarán sucesivamente su peso económico, sus efectos negativos y las alternativas propuestas.

En primer lugar, el turismo sigue siendo un motor esencial de la economía. La patronal Exceltur (documento 2) preveía que la actividad, impulsada por el regreso de los turistas asiáticos, creciera un 8,6 \% en 2024 y aportara por primera vez más de 200.000 millones de euros. Además, el documento 3 recuerda que, tras la crisis de 2008, el capital encontró en el turismo una nueva fuente de beneficios, y que la pandemia reveló hasta qué punto muchos territorios dependen de una sola actividad.

Sin embargo, este éxito tiene un coste social y ambiental cada vez más visible. El documento 1, firmado por plataformas vecinales, denuncia una industria descontrolada que expulsa a los habitantes de los centros históricos de Palma, Cádiz o Barcelona y convierte el patrimonio común en el negocio de unos pocos. Pese al turismo, Cádiz mantiene un paro superior al 20 \% y el barrio de La Viña ha perdido la quinta parte de su población en diez años. A ello se suman el consumo de agua en un país amenazado por la sequía y la destrucción de la posidonia en Baleares. La viñeta del documento 5 resume este proceso con humor parodiando «13, Rue del Percebe»: todos los pisos se han convertido en alojamientos de Airbnb o Booking, la planta baja alberga un Starbucks y al residente solo le queda una buhardilla de 20 metros cuadrados por 700 euros. Incluso las grandes empresas del documento 2 atribuyen el rechazo ciudadano al auge de los pisos turísticos.

Por ello, el modelo se cuestiona en la calle. En abril de 2024, las manifestaciones de «Canarias tiene un límite» abrieron un nuevo ciclo de protestas que une a ecologistas, vecinos y sindicatos (documento 3). La viñeta de Padylla (documento 4) rechaza la etiqueta de «turismofobia»: a unos turistas que preguntan, en un español torpe, si los manifestantes son «turismofóbicos», estos responden que son «canariasfílicos», es decir, que defienden su tierra sin odiar al visitante. Sus pancartas reclaman un turismo sostenible, una moratoria y una ecotasa. El documento 3 va más lejos y propone decrecer y transformar: ningún aeropuerto nuevo, el fin de la promoción internacional y programas públicos de turismo social, aunque advierte del riesgo de reservar los viajes a las élites. El documento 1, por su parte, exige límites urgentes y un diálogo con la ciudadanía.

En definitiva, el modelo turístico español evoluciona del crecimiento ilimitado hacia una exigencia de sostenibilidad: vecinos, activistas e incluso parte del sector reconocen que las ciudades deben seguir siendo, ante todo, lugares para vivir.

\textbf{(490 palabras, título incluido)}

Choix justifiés :

\begin{itemize}
\tightlist
\item
  La problématique est reformulée (\emph{está pasando de la expansión sin freno a la búsqueda de límites}) au lieu d'être recopiée, ce que le jury a bonifié.
\item
  \emph{preveía que la actividad \ldots{} creciera \ldots{} y aportara} : subjonctif imparfait après un verbe au passé (concordance des temps), structure que le jury dit rare et attendue.
\item
  Article devant les pourcentages (\emph{un 8,6 \%}, \emph{al 20 \%}), \emph{a} devant le complément de personne (\emph{expulsa a los habitantes}, \emph{une a ecologistas}), accent sur \emph{qué} interrogatif (\emph{en qué medida}).
\item
  La conclusion ne se termine pas par une question et n'exprime aucun avis.
\end{itemize}

\begin{methode}

Analyser un dessin humoristique en trois temps : ce que l'on voit (personnages, textes), le procédé (jeu de mots, parodie, exagération), la critique qu'il porte et le document qu'il rejoint.

\end{methode}

\begin{piege}

Pas de tutoiement ni d'adresse au lecteur, pas de \emph{yo pienso}. Attention aux diphtongues (\emph{demuestra}, \emph{encuentra}), à \emph{ser} / \emph{estar}, à \emph{hecho} avec un h, et au prétérit (\emph{abrieron}, \emph{reveló}) pour des faits datés.

\end{piege}

\begin{jury}

Rapport 2026 : « Les documents 4 et 5, notamment, ont souvent été traités de manière trop rapide, alors que leur dimension humoristique constituait un élément essentiel de leur portée critique. » (p.~13). « La synthèse n'est pas un résumé de chaque document mais une réflexion à construire à partir de chacun d'eux. » et « Nous rappelons que le tutoiement est à proscrire dans ce genre de synthèse » (p.~14). Pour le titre, le jury recommande « de privilégier des titres courts, dynamiques et percutants » (p.~14). Pour la conclusion, « il est vivement déconseillé de terminer par une question » (p.~15).

\end{jury}

\hypertarget{partie-v-italien}{%
\section{Partie V --- Italien}\label{partie-v-italien}}

Dossier sur le surtourisme en Italie : ses effets sur les résidents et l'environnement, les mesures locales (ticket d'entrée, quotas, réservation) et leurs limites.

\hypertarget{question-5-italien-synthuxe8se}{%
\subsection{Question 5 --- Italien, synthèse}\label{question-5-italien-synthuxe8se}}

\begin{indication}{1}

La question a deux volets (\emph{subisce} et \emph{fa fronte}) : votre plan doit les traiter tous deux, à parts à peu près égales.

\end{indication}

\begin{indication}{2}

Le tableau du document 4 ne se recopie pas : dégagez-en une tendance (concentration des nuitées sur quelques régions et quelques pays d'origine) et un ou deux ordres de grandeur.

\end{indication}

\begin{indication}{3}

La photo de la fontaine de Trevi doit être interprétée : que disent la foule, la saison, l'état de la visite ? Et le sondage du document 2 permet de nuancer l'efficacité des mesures.

\end{indication}

\hypertarget{solution-4}{%
\subsubsection{Solution}\label{solution-4}}

\textbf{Analyse du dossier.}

\begin{longtable}[]{@{}
  >{\raggedright\arraybackslash}p{(\columnwidth - 6\tabcolsep) * \real{0.0753}}
  >{\raggedright\arraybackslash}p{(\columnwidth - 6\tabcolsep) * \real{0.2322}}
  >{\raggedright\arraybackslash}p{(\columnwidth - 6\tabcolsep) * \real{0.5026}}
  >{\raggedright\arraybackslash}p{(\columnwidth - 6\tabcolsep) * \real{0.1598}}@{}}
\toprule\noalign{}
\begin{minipage}[b]{\linewidth}\raggedright
Doc.
\end{minipage} & \begin{minipage}[b]{\linewidth}\raggedright
Nature
\end{minipage} & \begin{minipage}[b]{\linewidth}\raggedright
Idée principale
\end{minipage} & \begin{minipage}[b]{\linewidth}\raggedright
Place dans le plan
\end{minipage} \\
\midrule\noalign{}
\endhead
\bottomrule\noalign{}
\endlastfoot
1 & article (\emph{leomagazineofficial.it}, 2024) & définition de l'OMT ; loyers, précarité, départ des résidents (Naples, Florence) ; ticket d'entrée de 5 € à Venise ; locations courtes de nouveau autorisées à Florence & I, II \\
2 & article (ANSA, 2024) & 49 \% des Italiens favorables aux mesures de contrôle, 38,4 \% contre ; hausse des taxes de séjour efficace pour les recettes, inefficace contre les arrivées & II \\
3 & article (\emph{italiaatavola.net}, 2024) & Cinq Terres (3500 habitants, plus de 3 millions de visiteurs), Dolomites, lac de Tenno ; énergie, croisières, déchets ; équilibre économie / territoire, tourisme œnogastronomique & I, II, conclusion \\
4 & tableau (Federalberghi / Istat, 2022) & 93,3 millions de nuitées européennes, dont 32,3 \% d'Allemands ; Trentin-Haut-Adige 22,5 \%, puis Vénétie et Lombardie & I \\
5 & photo (fontaine de Trevi) & foule estivale compacte : visite dégradée, monument exposé & I \\
\end{longtable}

\textbf{Plan.} I. Un pays qui subit un tourisme concentré (docs 4, 1, 3, 5). II. Des réponses locales, contestées et souvent insuffisantes (docs 1, 3, 2). Conclusion : concilier économie et protection du territoire.

\textbf{Sovraturismo: il prezzo del successo italiano}

Il dossier riunisce tre articoli del 2024, una tabella statistica e una fotografia. Tutti trattano dell'overtourism, che l'Organizzazione mondiale del turismo definisce come l'impatto eccessivo e negativo dei visitatori sulla qualità della vita dei residenti. Ci si chiede quindi in che misura l'Italia ne subisca le conseguenze e quali risposte riesca a darvi. Si esamineranno prima i danni del sovraffollamento, poi le misure adottate e i loro limiti.

Innanzitutto, l'Italia subisce gli effetti di un successo turistico concentrato in pochi luoghi. Come mostra la tabella del documento 4, nel 2022 i turisti europei hanno trascorso oltre 93 milioni di notti negli alberghi italiani: quasi un terzo di queste presenze era tedesco e più di un quinto si concentrava nel solo Trentino-Alto Adige, davanti al Veneto e alla Lombardia. Nelle città, l'afflusso fa salire gli affitti, alimenta il lavoro precario e spinge i residenti fuori dal centro: a Firenze gli affitti sono aumentati di oltre il 20\% da giugno 2023, mentre a Napoli chi vuole restare deve spesso accettare un impiego precario nella ristorazione (documento 1). Il documento 3 cita le Cinque Terre, dove 3.500 abitanti accolgono più di tre milioni di visitatori all'anno, e le Dolomiti, il cui ecosistema è messo a dura prova. A ciò si aggiungono l'inquinamento delle navi da crociera e i rifiuti: a Venezia un cestino si riempie in un quarto d'ora. La fotografia del documento 5 illustra questa saturazione: in piena estate, una folla compatta occupa tutta la piazza davanti alla fontana di Trevi, rendendo la visita scomoda e rischiosa per il monumento stesso.

Per far fronte al fenomeno, le amministrazioni sperimentano diverse soluzioni. Venezia ha sperimentato, per 29 giornate, un contributo d'accesso di 5 euro per i turisti che non pernottano in città, mentre le Cinque Terre hanno limitato gli ingressi e imposto la prenotazione online. Nelle Dolomiti si punta su sensori per monitorare i flussi, parcheggi di interscambio e mobilità sostenibile, mentre il lago di Tenno è presentato come un modello di gestione (documento 3). A Firenze, invece, la delibera che bloccava i nuovi affitti brevi in centro è caduta, e la nuova sindaca chiede un intervento nazionale. L'opinione pubblica appare divisa: secondo il sondaggio Jfc citato nel documento 2, il 49\% degli italiani approva misure di controllo dei flussi, ma il 38,4\% vi si oppone, in parte in nome della libertà di circolare nei luoghi pubblici. Infine, molte misure restano insufficienti: per Massimo Feruzzi, aumentare la tassa di soggiorno fa crescere le entrate dei comuni senza ridurre gli arrivi.

In conclusione, l'Italia subisce pienamente il sovraffollamento, nelle città d'arte come nei luoghi naturali. Le risposte esistono, ma restano locali e spesso parziali. Come suggerisce il documento 3, la sfida consiste nel conciliare i benefici economici del turismo con la tutela del territorio, per esempio attraverso un turismo enogastronomico capace di valorizzare luoghi meno conosciuti.

\textbf{(478 parole, titolo compreso)}

Choix justifiés :

\begin{itemize}
\tightlist
\item
  Le titre contient l'idée de surtourisme (\emph{Sovraturismo}), exigence explicite du jury.
\item
  \emph{Ci si chiede \ldots{} in che misura l'Italia ne subisca \ldots{} e quali risposte riesca a darvi} : interrogative indirecte au subjonctif.
\item
  \emph{rendendo la visita scomoda} : le gérondif a pour sujet celui de la principale (\emph{una folla compatta}), comme l'exige la grammaire italienne.
\item
  \emph{Infine} introduit le dernier argument, et non \emph{finalmente}.
\end{itemize}

\begin{methode}

Un tableau statistique se synthétise en une phrase de tendance et un ou deux chiffres arrondis (\emph{oltre 93 milioni}, \emph{quasi un terzo}, \emph{più di un quinto}), jamais en liste de valeurs.

\end{methode}

\begin{piege}

\emph{a} (préposition) et \emph{ha} (verbe \emph{avere}) ; \emph{finalmente} (enfin, après une attente) et \emph{infine} (pour finir) ; gérondif dont le sujet n'est pas celui de la principale (il faut alors une relative). Attention aussi au document 5 : le rapport lui-même l'appelle « doc. 1 » par erreur, c'est bien le document 5.

\end{piege}

\begin{jury}

Rapport 2026 (p.~16) : « Or il était indispensable de faire apparaître dans le titre l'idée de surtourisme (sovra-turismo o affollamento turistico). » Le document iconographique « a été traité de manière superficielle, quand il n'a pas été totalement escamoté », et « il ne suffit pas de citer vaguement « la situation que l'on voit au doc. 5 » pour considérer avoir traité le document. » Sur la langue : « finalmente ne correspond pas à finalement » et « le gérondif doit être référé au sujet ».

\end{jury}

\hypertarget{bilan}{%
\section{Bilan}\label{bilan}}

\textbf{À retenir pour toutes les langues.}

\begin{itemize}
\tightlist
\item
  Lire tout le dossier avant d'écrire, dresser une grille (document, idée principale, chiffres, ton), puis construire un plan qui croise les documents au lieu de les résumer l'un après l'autre.
\item
  Introduction courte : une accroche tirée du dossier et reformulée, la problématique reformulée, le plan annoncé. Développement en deux ou trois parties équilibrées. Conclusion qui répond à la question, sans avis personnel ni question finale.
\item
  Citer chaque document explicitement, images et tableaux compris, et les interpréter.
\item
  Un titre court, précis, qui contient le thème du dossier ; le nombre de mots en fin de copie, et des repères tous les 50 mots.
\item
  Viser 470 à 490 mots : la marge protège des écarts de convention de décompte (élisions, nombres).
\end{itemize}

\textbf{Ce que les jurys 2026 ont jugé discriminant.} En anglais, la restitution exacte des coûts, de la position de Spinrad et de l'ironie du dessin, et la mention des \emph{space mirrors} ; en allemand, l'analyse de la photo et l'équilibre entre les documents ; en espagnol, l'analyse de l'humour des dessins et la reformulation de la problématique ; en italien, le surtourisme dans le titre et l'interprétation de la photo ; en arabe, une problématique bien formulée et la correction de la langue.

\textbf{Grammaire à revoir, langue par langue.} Allemand : place du verbe dans les subordonnées, déclinaisons, genres, virgules. Anglais : \emph{present perfect} avec \emph{since}, interrogative indirecte, indénombrables (\emph{research}, \emph{damage}), pas d'article devant \emph{climate change}. Arabe : hamza de liaison, prépositions régies par les verbes, accords. Espagnol : subjonctif (\emph{para que}, concordance), article devant \%, diphtongues, \emph{a} devant le complément de personne. Italien : \emph{a} / \emph{ha}, gérondif, faux amis (\emph{finalmente}).

\textbf{Situations qui se retrouvent dans d'autres sujets.} Un progrès technique présenté comme solution miracle (géo-ingénierie, intelligence artificielle) ; une prospérité économique qui se retourne contre les habitants (tourisme de masse) ; une loi qui existe mais s'applique mal (discriminations) ; une ressource vitale sous tension (eau). Les mêmes plans fonctionnent : constat, réponses, limites.

\end{document}
